\section*{Chapitre \ref{ch:b3:hilbert} --- Espaces de Hilbert}

\begin{solution}{exo:b3:hilbert:1}
(a) Réel~: développer $\norm{x \pm y}^2 = \norm x^2 \pm 2\langle
x,y\rangle + \norm y^2$ et soustraire. Complexe (\omterm{def:b3:hilbert:inner}{produit
scalaire} linéaire en le second argument)~: en développant comme
ci-dessus,
\[
\langle x, y\rangle = \frac14\sum_{k=0}^{3}
\iu^k\,\bigl\|\iu^kx + y\bigr\|^2,
\]
chaque terme contribuant $\iu^k\cdot2\operatorname{Re}\bigl(
(-\iu)^k\langle x,y\rangle\bigr)$, dont la somme est $4\langle
x,y\rangle$ (vérifier les quatre valeurs de $k$~; le
$\sum\iu^k (\norm x^2 + \norm y^2) = 0$).

(b) $L^1$~: $f = \mathbf 1_{\intcc0{1/2}}$, $g = \mathbf
1_{\intcc{1/2}1}$~: $\norm{f\pm g}_1^2 = 1$ chacun, somme $2$~;
$2\norm f_1^2 + 2\norm g_1^2 = 1 \neq 2$. Norme uniforme~: $f =
\mathbf 1$, $g(t) = t$ sur $\intcc01$~: $\norm{f + g}_\infty^2 +
\norm{f-g}_\infty^2 = 4 + 1 = 5 \neq 4 = 2 + 2$. En échouant
l'identité du parallélogramme, ces normes ne sont induites par
aucun \omterm{def:b3:hilbert:inner}{produit scalaire} (qui la forcerait par développement
direct).
\end{solution}

\begin{solution}{exo:b3:hilbert:2}
(a) $p(f) = \bigl(\int_0^1f\bigr)\mathbf 1$~: en effet $f - \int f
\perp$ les constantes ($\int(f - \int f)c = 0$). La meilleure
approximation constante de $f$ en moyenne quadratique est sa
\emph{moyenne} --- la première instance d'espérance conditionnelle
(\cref{ch:b3:probability}).

(b) $p(f) = f\,\mathbf 1_{\intcc{1/2}1}$~: la différence
$f\mathbf 1_{\intcc0{1/2}}$ est orthogonale à tout $g$
s'annulant sur $\intcc0{1/2}$.

(c) $d^2 = \bigl\|x - \tfrac12\bigr\|_2^2 = \int_0^1(x -
\tfrac12)^2\dd x = \tfrac1{12}$~: $d = \frac1{2\sqrt3}$.
\end{solution}

\begin{solution}{exo:b3:hilbert:3}
(a) L'équivalence est \cref{thm:b3:hilbert:decomposition}
($\bar F = (F^\perp)^\perp$, et $\bar F = H \iff F^\perp =
\{0\}$). Exemple~: l'espace $F$ des suites finies est \omterm{def:b3:topology:interior}{dense}
dans $\ell^2$ (troncature) et propre~: $F^\perp = \{0\}$
pourtant $F \neq \ell^2$ --- pour un sous-espace non fermé, $H
= F \oplus F^\perp$ échoue ostensiblement ($F \oplus \{0\}
\neq H$).

(b) $\abs{\langle x_n, y_n\rangle - \langle x, y\rangle} \leq
\abs{\langle x_n - x, y_n\rangle} + \abs{\langle x, y_n -
y\rangle} \leq \norm{x_n - x}\sup_n\norm{y_n} + \norm
x\,\norm{y_n - y} \to 0$ (les suites convergentes sont
bornées). Utilisé~: dans \cref{thm:b3:hilbert:parseval}(1)
pour voir $x - \lim S_N \perp e_k$, et dans
\cref{thm:b3:hilbert:decomposition} pour voir $F^\perp =
\bar F^{\,\perp}$.
\end{solution}

\begin{solution}{exo:b3:hilbert:4}
$e_0 = \frac1{\sqrt2}$. Ensuite, $x \perp \mathbf 1$ déjà
($\int_{-1}^1x = 0$), et $\int_{-1}^1x^2 = \frac23$~: $e_1 =
\sqrt{\tfrac32}\,x$. Puis $x^2 - \langle e_0, x^2\rangle e_0 =
x^2 - \frac13$ (et $\perp e_1$ par parité), avec
\[
\int_{-1}^1\Bigl(x^2 - \frac13\Bigr)^2\dd x = \frac25 -
\frac49 + \frac29 = \frac{8}{45}:
\qquad e_2 = \sqrt{\tfrac{45}8}\,\Bigl(x^2 - \frac13\Bigr).
\]
Comparaison~: $P_0 = 1$, $P_1 = x$, $P_2 = \frac{3x^2 - 1}2$, et
$\sqrt{n + \tfrac12}\,P_n$ donne $\frac1{\sqrt2}$,
$\sqrt{\frac32}x$, $\sqrt{\frac52}\,\frac{3x^2-1}2 =
\sqrt{\frac{45}8}\bigl(x^2 - \frac13\bigr)$~: exactement $e_0,
e_1, e_2$.
\end{solution}

\begin{solution}{exo:b3:hilbert:5}
Pour $f(t) = t$~: $c_0 = 0$ et, en intégrant par parties, $c_n =
\frac{\iu(-1)^n}{n}$ pour $n \neq 0$~: $\abs{c_n}^2 =
\frac1{n^2}$. \omterm{thm:b3:hilbert:parseval}{Parseval}~:
\[
\frac1{2\pi}\int_{-\pi}^\pi t^2\dd t = \frac{\pi^2}3
= \sum_{n\neq0}\frac1{n^2} = 2\sum_{n\geq1}\frac1{n^2}
\ \Longrightarrow\ \sum_{n\geq1}\frac1{n^2} = \frac{\pi^2}6 .
\]
Pour $f(t) = t^2$~: $c_0 = \frac{\pi^2}3$, $c_n =
\frac{2(-1)^n}{n^2}$ ($n \ne 0$). \omterm{thm:b3:hilbert:parseval}{Parseval}~:
\[
\frac1{2\pi}\int_{-\pi}^{\pi}t^4\dd t = \frac{\pi^4}5
= \frac{\pi^4}9 + \sum_{n\neq0}\frac4{n^4}
\ \Longrightarrow\
\sum_{n\geq1}\frac1{n^4} = \frac18\Bigl(\frac{\pi^4}5 -
\frac{\pi^4}9\Bigr) = \frac{\pi^4}{90} .
\]
Aucune précaution $\mathcal C^1$ par morceaux n'est nécessaire~:
\cref{thm:b3:hilbert:fourier} couvre toute fonction $L^2$.
\end{solution}

\begin{solution}{exo:b3:hilbert:6}
(a) $\varphi(f) = \int_0^{1/2}f = \langle\mathbf
1_{\intcc0{1/2}},\ f\rangle$~: le représentant est $a =
\mathbf 1_{\intcc0{1/2}}$, et $\norm\varphi = \norm a_2 =
\frac1{\sqrt2}$ (\cref{thm:b3:hilbert:riesz}).

(b) Prendre les fonctions tente $f_n$ de pic $1$ en
$\frac12$ et support de largeur $\frac2n$~: $f_n(\tfrac12) = 1$
tandis que $\norm{f_n}_2^2 \leq \frac2n \to 0$~: aucune
constante $C$ ne peut donner $\abs{f(\frac12)} \leq C\norm
f_2$. L'évaluation ponctuelle est sans sens dans $L^2$ ---
les éléments sont des classes modulo les ensembles nuls ---
et ce calcul en est la raison quantitative.
\end{solution}

\begin{solution}{exo:b3:hilbert:7}
(a) Soit $M = \sup_k\norm{x_k}$. Les suites scalaires
$(\langle e_n, x_k\rangle)_k$ sont bornées par $M$~: une
extraction diagonale produit $x_{k_j}$ avec $\langle e_n,
x_{k_j}\rangle \to \gamma_n$ pour tout $n$. Pour chaque $N$~:
$\sum_{n\leq N}\abs{\gamma_n}^2 = \lim_j\sum_{n\leq
N}\abs{\langle e_n, x_{k_j}\rangle}^2 \leq M^2$ (Bessel), donc
$(\gamma_n) \in \ell^2$ et $x = \sum_n\gamma_ne_n \in H$
(\cref{thm:b3:hilbert:parseval}(3)). Pour $y \in H$~:
\[
\abs{\langle y, x_{k_j} - x\rangle}
\leq \Bigl|\sum_{n\leq N}\overline{c_n(y)}\bigl(\langle e_n,
x_{k_j}\rangle - \gamma_n\bigr)\Bigr|
+ 2M\Bigl(\sum_{n>N}\abs{c_n(y)}^2\Bigr)^{1/2},
\]
en utilisant le développement $\langle y, z\rangle =
\sum\overline{c_n(y)}c_n(z)$ et Cauchy--Schwarz sur la queue~;
choisir $N$ puis $j$~: convergence faible vers $x$.

(b) $\langle y, e_n\rangle = c_n(y) \to 0$ pour tout $y$
(queues $\ell^2$)~: $e_n \rightharpoonup 0$, pourtant
$\norm{e_n} = 1$~: la norme n'est pas faiblement \omterm{def:b3:topology:continuity}{continue}. En
(a)~: $\norm x^2 = \sum\abs{\gamma_n}^2 \leq
\liminf_j\norm{x_{k_j}}^2$ (sections finies et Bessel encore)~:
les limites faibles ne peuvent que perdre de la norme.
\end{solution}

\begin{solution}{exo:b3:hilbert:8}
Pour $y$ fixé, $x \mapsto \langle y, Tx\rangle$ est une forme
linéaire \omterm{def:b3:topology:continuity}{continue}~; Riesz donne un unique $T^*y$ avec
$\langle y, Tx\rangle = \langle T^*y, x\rangle$ pour tout $x$
--- en conjuguant, $\langle Tx, y\rangle = \langle x,
T^*y\rangle$. L'unicité rend $T^*$ linéaire~;
\[
\norm{T^*y} = \sup_{\norm x = 1}\abs{\langle T^*y, x\rangle}
= \sup_{\norm x=1}\abs{\langle y, Tx\rangle}
\leq \vertiii T\,\norm y,
\]
donc $\vertiii{T^*} \leq \vertiii T$, et $T^{**} = T$ donne
l'égalité. Décalage~: $\langle Sx, y\rangle = \sum_{n\geq1}
x_n\bar y_{n+1} = \langle x, S^*y\rangle$ avec $(S^*y)_n =
y_{n+1}$~: le décalage rétrograde. Noyau--image~: $T^*y = 0$
ssi $\langle x, T^*y\rangle = 0$ pour tout $x$ ssi $\langle
Tx, y\rangle = 0$ pour tout $x$ ssi $y \perp
\operatorname{im}T$~: $\ker T^* = (\operatorname{im}T)^\perp$~;
en prenant $\perp$ et en utilisant
le~\cref{thm:b3:hilbert:decomposition},
$\overline{\operatorname{im}T} = (\ker T^*)^\perp$.
\end{solution}

\begin{solution}{exo:b3:hilbert:9}
Si $a(u, \cdot) = \varphi$~: pour tout $w$,
\[
J(u + w) - J(u) = a(u, w) - \varphi(w) + \tfrac12a(w,w)
= \tfrac12a(w,w) \geq \tfrac\alpha2\norm w^2,
\]
strictement positif pour $w \neq 0$~: $u$ est l'unique
minimiseur. Réciproquement, en un minimiseur la fonction $t
\mapsto J(u + tw)$ (un polynôme quadratique en $t$) a une
dérivée nulle en $0$~: $a(u, w) - \varphi(w) = 0$ pour tout
$w$. Projection redérivée~: pour un sous-espace fermé $F$,
appliquer Lax--Milgram sur l'\omterm{def:b3:hilbert:inner}{espace de Hilbert} $F$ avec
$a(u,v) = \langle u, v\rangle$ ($M = \alpha = 1$) et
$\varphi(v) = \langle x, v\rangle$~: un unique $p \in F$ avec
$\langle p, v\rangle = \langle x, v\rangle$ pour tout $v \in
F$, i.e.\ $x - p \perp F$ --- et par le cas symétrique, $p$
minimise $\frac12\norm v^2 - \langle x, v\rangle =
\frac12\norm{v - x}^2 - \frac12\norm x^2$ sur $F$~: la
projection.
\end{solution}

\begin{solution}{exo:b3:hilbert:10}
Normalisation~: $\int h_n^2 = 2^j\cdot 2^{-j} = 1$.
Orthogonalité~: deux fonctions de Haar distinctes ont soit des
supports (intérieurs) disjoints (produit nul p.p.), soit le
support de la plus fine \omterm{def:b3:rings:content}{contenu} dans un demi-intervalle où la
plus grossière est constante --- alors l'intégrale du produit
est cette constante fois $\int h_{\text{fine}} = 0$~; contre
$h_0 = \mathbf 1$, encore $\int h_n = 0$. Totalité~: l'enveloppe
de $\{h_0, \dots, h_{2^J-1}\}$ est constituée des fonctions en
escalier sur la grille dyadique de pas $2^{-J}$~; les deux
espaces ont dimension $2^J$ et les fonctions de Haar sont
indépendantes (orthonormées)~: l'enveloppe est \emph{toutes}
ces fonctions en escalier. Les fonctions en escalier dyadiques
sont \omterm{def:b3:topology:interior}{denses} dans $L^2(\intcc01)$~: les \omterm{def:b3:lebesgue:simple}{fonctions simples} sont
\omterm{def:b3:topology:interior}{denses} (\cref{thm:b3:lp:density}(1)), les ensembles
\omterm{def:b3:lebesgue:measurable}{mesurables} s'approchent par unions finies d'intervalles
(\cref{exo:b3:measure:7}), et les intervalles par des
dyadiques (extrémités bougent de $\leq 2^{-J}$). Par
le~\cref{thm:b3:hilbert:parseval}, le système de Haar est une
\omterm{def:b3:hilbert:onb}{base hilbertienne}.
\end{solution}

\begin{solution}{exo:b3:hilbert:11}
(i) $\Rightarrow$ (ii)~: pour la projection orthogonale,
$\langle Px, y\rangle = \langle Px, Py\rangle = \langle x,
Py\rangle$ (insérer les décompositions $x = Px + (x - Px)$
etc.\ et tuer les termes croisés). (ii) $\Rightarrow$ (iii)~:
$\norm{Px}^2 = \langle P^2x, x\rangle = \langle Px, x\rangle
\leq \norm{Px}\norm x$, donc $\vertiii P \leq 1$, et $Pu = u$
sur l'image non nulle~: $= 1$. (iii) $\Rightarrow$ (i)~: $H =
\operatorname{im}P \oplus \ker P$ (algébriquement, de $P^2
= P$)~; prendre $u = Pu \in \operatorname{im}P$, $v \in \ker
P$, $t \in \R$~: $\norm{P(u + tv)}^2 = \norm u^2$ doit être
$\leq \norm{u + tv}^2 = \norm u^2 + 2t\operatorname{Re}\langle
u, v\rangle + t^2\norm v^2$ pour tout $t$, forçant
$\operatorname{Re}\langle u, v\rangle = 0$ (comparer les
termes linéaires quand $t \to 0^\pm$)~; remplacer $v$ par
$\iu v$ tue aussi la partie imaginaire~:
$\operatorname{im}P \perp \ker P$, ce qui est exactement
l'orthogonalité de la projection. Exemple~: $P(x, y) = (x +
y, 0)$ sur $\R^2$~: $P^2 = P$, image l'axe des $x$, noyau la
droite $y = -x$, et $\vertiii P =
\sup\frac{\abs{x+y}}{\norm{(x,y)}} = \sqrt2$ (atteint en
$(1,1)/\sqrt2$)~: une projection oblique a une norme $> 1$.
(Pour la chronique, (ii) donne aussi (i) directement~: $\ker P
= \ker P^* = (\operatorname{im}P)^\perp$ par
l'~\cref{exo:b3:hilbert:8}.)
\end{solution}

\begin{solution}{exo:b3:hilbert:12}
(a) Pour $U$ unitaire~: $\norm{Ux - x}^2 = 2\norm x^2 -
2\operatorname{Re}\langle Ux, x\rangle$ et $\norm{U^*x -
x}^2 = 2\norm x^2 - 2\operatorname{Re}\langle x, Ux\rangle$~:
les deux s'annulent ensemble, donc $\ker(U - I) = \ker(U^* -
I)$. Puis, en utilisant $\ker T^* = (\operatorname{im}T)^\perp$
(\cref{exo:b3:hilbert:8}) avec $T = U - I$ et $T^* = U^* -
I$~:
\[
\overline{\operatorname{im}(U - I)} = \bigl(\ker(U^* -
I)\bigr)^\perp = F^\perp .
\]

(b) Sur $F$~: $U^kx = x$, donc $A_nx = x$. Pour $x = (U -
I)y$~: $A_nx = \frac1n(U^ny - y)$, de norme $\leq
\frac2n\norm y \to 0$. Pour $x$ dans la fermeture~: pour
$\varepsilon$ donné, prendre $x' = (U - I)y$ avec $\norm{x -
x'} < \varepsilon$~; comme $\vertiii{A_n} \leq
\frac1n\sum\vertiii{U^k} = 1$, $\norm{A_nx} \leq
\norm{A_n(x - x')} + \norm{A_nx'} \leq \varepsilon + o(1)$.

(c) Décomposer $x = Px + (x - Px)$ avec $Px \in F$ et $x -
Px \in F^\perp = \overline{\operatorname{im}(U - I)}$
(partie (a))~: $A_nx = Px + A_n(x - Px) \to Px + 0$.

(d) Dans la \omterm{def:b3:topology:basis}{base} de Fourier $e_m(x) = \eu^{2\iu\pi mx}$~:
$Ue_m = \eu^{2\iu\pi m\alpha}e_m$, donc $Ue_m = e_m$ ssi
$m\alpha \in \Z$ ssi $m = 0$ ($\alpha$ irrationnel)~: $F =
\C\mathbf 1$ et $Pf = \langle\mathbf 1, f\rangle\mathbf 1 =
\int_0^1f$. Le théorème se lit $\frac1n\sum_{k<n}f(\cdot +
k\alpha) \to \int_0^1f$ dans $L^2(\R/\Z)$~: les moyennes
d'\omterm{def:b3:groups:action}{orbite} d'une rotation irrationnelle s'équidistribuent ---
l'ombre $L^2$ du théorème d'équidistribution de Weyl, obtenue
par pure géométrie hilbertienne.
\end{solution}

\begin{solution}{pb:b3:hilbert:1}
\textbf{1.} Gram--Schmidt garantit
$\operatorname{Vect}(p_0, \dots, p_n) =
\operatorname{Vect}(1, \dots, t^n) = \R_n[t]$ et $p_n \perp
p_k$ ($k < n$), d'où $p_n \perp \R_{n-1}[t]$. Les $p_k$, de
\omterm{def:b3:galois:extension}{degrés} strictement croissants, sont indépendants~: une \omterm{def:b3:topology:basis}{base}.

\textbf{2.} $t\,p_n$ est monique de \omterm{def:b3:galois:extension}{degré} $n + 1$~: développer
$t\,p_n = p_{n+1} + \sum_{k\leq n}c_kp_k$ avec $c_k =
\langle p_k, tp_n\rangle/\norm{p_k}^2$. Pour $k \leq n - 2$~:
$\langle p_k, tp_n\rangle = \langle tp_k, p_n\rangle = 0$
(\omterm{def:b3:galois:extension}{degré} $k + 1 < n$). Donc $tp_n = p_{n+1} + a_np_n +
b_np_{n-1}$, la récurrence annoncée, avec
\[
b_n = \frac{\langle p_{n-1}, tp_n\rangle}{\norm{p_{n-1}}^2}
= \frac{\langle tp_{n-1}, p_n\rangle}{\norm{p_{n-1}}^2}
= \frac{\langle p_n + (\text{inférieurs}),\
p_n\rangle}{\norm{p_{n-1}}^2}
= \frac{\norm{p_n}^2}{\norm{p_{n-1}}^2} > 0 .
\]

\textbf{3.} Soit $t_1 < \dots < t_m$ les points intérieurs à
$I$ où $p_n$ change de signe, et $q = \prod_{i\leq m}(t -
t_i)$ (avec $q = 1$ si $m = 0$). Alors $p_nq$ a un signe
constant sur $I$ et n'est pas p.p.\ nul~: $\int_Ip_nq\,w \neq
0$. Si $m < n$, ceci contredit $p_n \perp \R_{n-1}[t]$. Donc
$m = n$~: $p_n$ a $n$ racines intérieures distinctes (il en a
au plus $n$ en tout).

\textbf{4.} $(t^2 - 1)^n$ a \omterm{def:b3:galois:extension}{degré} $2n$~; $n$ dérivées laissent
\omterm{def:b3:galois:extension}{degré} $n$, avec coefficient dominant
$\frac{(2n)(2n-1)\cdots(n+1)}{2^nn!} =
\frac{(2n)!}{2^n(n!)^2}$. Pour $\deg Q < n$, intégrer par
parties $n$ fois~: tous les termes de bord contiennent une
dérivée d'ordre $< n$ de $(t^2-1)^n$, qui s'annule en $\pm1$
(racine d'ordre $n$)~; après $n$ étapes l'intégrande porte
$Q^{(n)} = 0$.

\textbf{5.} Avec $u = (t^2 - 1)^n$~:
\[
(2^nn!)^2\norm{P_n}^2 = \int_{-1}^1(u^{(n)})^2
= (-1)^n\int_{-1}^1 u\,u^{(2n)}
= (2n)!\int_{-1}^1(1 - t^2)^n\dd t ,
\]
($u^{(2n)} = (2n)!$~; termes de bord nuls comme en question
4). Et $\int_{-1}^1(1-t^2)^n\dd t = B(\tfrac12, n+1) =
\frac{\Gamma(\frac12)\Gamma(n+1)}{\Gamma(n + \frac32)} =
\frac{2\cdot4^n(n!)^2}{(2n+1)!}$
(\cref{exo:b3:product:8}). En combinant~: $\norm{P_n}^2 =
\frac{2}{2n + 1}$.

\textbf{6.} Les polynômes sont $\norm\cdot_\infty$-denses dans
$\mathcal C(\intcc{-1}1)$ (Weierstrass,
\cref{cor:b3:complete:weierstrass}), les fonctions \omterm{def:b3:topology:continuity}{continues}
sont $L^2$-denses (\cref{thm:b3:lp:density}), et
$\norm\cdot_2 \leq \sqrt2\norm\cdot_\infty$~: les enveloppes
polynomiales sont totales, donc les $P_n$ normalisés forment
une \omterm{def:b3:hilbert:onb}{base hilbertienne}. Développement de $\abs t$~: le
coefficient contre $P_0$ est $\frac{\langle P_0, \abs
t\rangle}{\norm{P_0}^2} = \frac12$~; contre $P_1$~: $0$
(parité)~; contre $P_2$~:
$\frac{\int_{-1}^1\abs t\,\frac{3t^2-1}2\dd t}{2/5} =
\frac{1/4}{2/5} = \frac58$. Meilleure approximation
quadratique~:
\[
\abs t \approx \frac12 + \frac58\,P_2(t) = \frac{3}{16} +
\frac{15}{16}\,t^2 .
\]

\textbf{7.} De $\frac{\dd^{n+1}}{\dd t^{n+1}}\eu^{-t^2} =
\frac{\dd^n}{\dd t^n}(-2t\,\eu^{-t^2})$ et Leibniz, $H_{n+1}
= 2tH_n - H_n'$~; la récurrence donne \omterm{def:b3:galois:extension}{degré} $n$ et
coefficient dominant $2^n$. Pour $m < n$, intégrer par
parties $n$ fois dans $\int H_m H_n\eu^{-t^2} =
(-1)^n\int H_m\,\bigl(\eu^{-t^2}\bigr)^{(n)}$~: les termes de
bord (polynôme $\times$ $\eu^{-t^2}$) s'annulent en
$\pm\infty$, laissant $\int H_m^{(n)}\,\eu^{-t^2} = 0$. Pour
$m = n$~: $H_n^{(n)} = 2^nn!$, donc $\norm{H_n}_w^2 =
2^nn!\int\eu^{-t^2} = 2^nn!\sqrt\pi$.

\textbf{8.} Soit $f \in L^2(\R, \eu^{-t^2}\dd t)$ orthogonal
à tout polynôme, et $g = f\eu^{-t^2}$. Alors $g \in L^1$~:
$\int\abs f\eu^{-t^2} \leq \bigl(\int\abs
f^2\eu^{-t^2}\bigr)^{1/2}\bigl(\int\eu^{-t^2}\bigr)^{1/2}$
(Cauchy--Schwarz). Pour $\xi \in \R$, développer
$\eu^{-\iu\xi t}$~: les sommes partielles sont dominées car
\[
\sum_k\frac{\abs\xi^k}{k!}\int\abs f\,\abs t^k\eu^{-t^2}\dd t
\leq \Bigl(\int \abs f^2\eu^{-t^2}\Bigr)^{1/2}
\sum_k\frac{\abs\xi^k}{k!}\Bigl(\int
t^{2k}\eu^{-t^2}\Bigr)^{1/2} < \infty
\]
(la dernière série converge~: $\int t^{2k}\eu^{-t^2} =
\Gamma(k+\frac12) \leq k!\,\sqrt\pi$, donc les termes sont
$O(\abs\xi^k/\sqrt{k!})$). L'intégration terme à terme
(\cref{cor:b3:lebesgue:additivity} appliqué à la série
absolue, puis Fubini pour les séries) donne
\[
\int_\R g(t)\,\eu^{-\iu\xi t}\dd t
= \sum_k\frac{(-\iu\xi)^k}{k!}\int f(t)\,t^k\,\eu^{-t^2}\dd t
= 0 ,
\]
chaque intégrale étant de type $\langle t^k, f\rangle_w = 0$.
Par l'injectivité admise de la transformée de Fourier
(\cref{ch:b3:fouriertransform}), $g = 0$ p.p., donc $f = 0$
p.p.~: la famille d'Hermite (dont les enveloppes sont les
polynômes) est totale.

\textbf{9.} Exactitude jusqu'au \omterm{def:b3:galois:extension}{degré} $n - 1$~: pour un tel
$P$, $P = \sum_iP(t_i)\ell_i$ exactement, donc $\int Pw =
\sum_iP(t_i)\int \ell_iw = Q(P)$. \omterm{def:b3:galois:extension}{Degré} $\leq 2n - 1$~:
diviser $P = qp_n + r$, $\deg q \leq n - 1$, $\deg r \leq
n-1$~; alors $\int Pw = \int qp_nw + \int rw = 0 + Q(r)$
($p_n \perp \R_{n-1}[t]$), tandis que $Q(P) =
\sum_iw_i\bigl(q(t_i)\,p_n(t_i) + r(t_i)\bigr) = Q(r)$
puisque les nœuds sont les racines de $p_n$. Égal.

\textbf{10.} $\ell_i^2$ a \omterm{def:b3:galois:extension}{degré} $2n - 2 \leq 2n - 1$ et
$\ell_i^2(t_j) = \delta_{ij}$~: $0 < \int\ell_i^2w =
Q(\ell_i^2) = w_i$. Polya (\cref{exo:b3:banach:9},
transporté à $I$ avec poids)~: la condition (i) tient ---
chaque polynôme est intégré exactement une fois $2n - 1
\geq$ son \omterm{def:b3:galois:extension}{degré}~; condition (ii)~: $\sum_i\abs{w_{i}} =
\sum_iw_i = Q(\mathbf 1) = \int_Iw$, bornée~: $Q_n(f) \to
\int fw$ pour toute $f \in \mathcal C(I)$, $I$ compact.

\textbf{11.} Monique $p_2 = t^2 - \frac13$ (de
l'~\cref{exo:b3:hilbert:4})~: nœuds $\pm\frac1{\sqrt3}$. Poids~:
$\ell_1(t) = \frac{t - \frac1{\sqrt3}}{-\frac2{\sqrt3}}$, et
$w_1 = \int_{-1}^1\ell_1 = 1$~; par symétrie $w_2 = 1$.
Exactitude~: $\int 1 = 2 = 1 + 1$~; $\int t = 0 =
-\frac1{\sqrt3} + \frac1{\sqrt3}$~; $\int t^2 = \frac23 =
\frac13 + \frac13$~; $\int t^3 = 0$. La règle des trapèzes à
deux points (nœuds $\pm1$, poids $1, 1$) n'est exacte que
jusqu'au \omterm{def:b3:galois:extension}{degré} $1$~: sur $t^2$ elle renvoie $2$ au lieu de
$\frac23$. Même coût, deux \omterm{def:b3:galois:extension}{degrés} d'exactitude en plus~: le
dividende des nœuds orthogonaux.

\textbf{12.} De $\cos(n{+}1)\theta + \cos(n{-}1)\theta =
2\cos\theta\cos n\theta$~: $T_{n+1} = 2tT_n - T_{n-1}$ avec
$T_0 = 1$, $T_1 = t$~; la récurrence donne des polynômes de
\omterm{def:b3:galois:extension}{degré} $n$ de coefficient dominant $2^{n-1}$ ($n \geq 1$).
En substituant $t = \cos\theta$ ($w(t)\dd t \mapsto
\dd\theta$)~: $\langle T_m, T_n\rangle_w =
\int_0^\pi\cos m\theta\cos n\theta\,\dd\theta = 0$ pour $m
\neq n$, $= \pi$ pour $m = n = 0$, $= \frac\pi2$ sinon
(produit-en-somme). \omterm{def:b3:galois:extension}{Degrés} et orthogonalité deux à deux
identifient les $T_n$ avec la sortie de Gram--Schmidt à
scalaires près~; un développement de Chebyshev de $f$ est
exactement la série de Fourier en cosinus de $\theta \mapsto
f(\cos\theta)$.

\textbf{13.} $T_n(t) = 0$ ssi $\cos n\theta = 0$ ssi
$\theta = \frac{(2k-1)\pi}{2n}$~: les $n$ racines distinctes
$t_k = \cos\frac{(2k-1)\pi}{2n} \in \intoo{-1}1$. Extrema~:
$\abs{T_n} \leq 1$ sur $\intcc{-1}1$, avec $T_n(s_j) =
(-1)^j$ aux $n + 1$ points $s_j = \cos\frac{j\pi}n$~:
équi-oscillation parfaite.

\textbf{14.} $2^{1-n}T_n$ est monique de norme uniforme
$2^{1-n}$. Si un monique $P$ de \omterm{def:b3:galois:extension}{degré} $n$ avait $\sup\abs P <
2^{1-n}$, la différence $D = 2^{1-n}T_n - P$ aurait \omterm{def:b3:galois:extension}{degré}
$\leq n - 1$ (termes dominants s'annulent) pourtant
alternerait en signe en $s_0 > \dots > s_n$ (là
$2^{1-n}T_n = \pm2^{1-n}$ domine $P$)~: au moins $n$ zéros
--- $D \equiv 0$, contradiction. Pour l'unicité à
l'égalité, le même $D$ satisfait $(-1)^jD(s_j) \geq 0$~; un
polynôme non nul de \omterm{def:b3:galois:extension}{degré} $\leq n-1$ ne peut avoir $n$
contraintes extrémales faiblement alternées sans $n$ racines
comptées correctement (si $D(s_j) = 0$ pour un $s_j$
intérieur, ce zéro est double dans le comptage car $D$ garde
un signe localement)~: encore $D \equiv 0$.

\textbf{15.} La formule d'erreur de Lagrange (Rolle, deuxième
année) donne $f - L_nf = \frac{f^{(n)}(\xi_t)}{n!}\,\omega(t)$,
donc l'erreur uniforme est au plus
$\frac{\norm{f^{(n)}}_\infty}{n!}\,\sup\abs\omega$, et
$\omega$ est monique de \omterm{def:b3:galois:extension}{degré} $n$~: par la question 14,
$\sup_{\intcc{-1}1}\abs\omega \geq 2^{1-n}$ avec égalité ssi
les nœuds sont les racines de Chebyshev. D'où la borne
optimale $\norm{f - L_nf}_\infty \leq
\frac{\norm{f^{(n)}}_\infty}{2^{n-1}n!}$. Avec des nœuds
équidistants, $\sup\abs\omega$ est exponentiellement plus
grand près des extrémités, et interpoler même $\frac1{1 +
25t^2}$ y diverge quand $n \to \infty$ (phénomène de Runge)~;
les nœuds de Chebyshev sont le remède.

\textbf{16.} En dérivant $T_n(\cos\theta) = \cos n\theta$~:
$T_n'(\cos\theta) = \frac{n\sin n\theta}{\sin\theta}$, qui
tend vers $n^2$ quand $\theta \to 0$ et vers
$(-1)^{n+1}n^2$ quand $\theta \to \pi$~:
$\abs{T_n'(\pm1)} = n^2$. Aux points intérieurs,
$\abs{T_n'(t)} \leq \frac{n}{\sqrt{1 - t^2}} = O(n)$~: le
souffle quadratique ne vit qu'aux bords (borne intérieure de
Bernstein versus borne globale de Markov).

\textbf{17.} Soit $\theta_k = \frac{(2k-1)\pi}{2n}$ et $S_j
= \sum_{k=1}^n\cos(j\theta_k)$ pour $1 \leq j \leq n-1$.
Alors
\[
S_j = \operatorname{Re}\Bigl[\eu^{\iu j\pi/2n}
\sum_{k=0}^{n-1}\eu^{\iu jk\pi/n}\Bigr]
= \operatorname{Re}\Bigl[\eu^{\iu j\pi/2n}\,
\frac{\eu^{\iu j\pi} - 1}{\eu^{\iu j\pi/n} - 1}\Bigr] .
\]
Pour $j$ pair le numérateur s'annule~: $S_j = 0$. Pour $j$
impair le numérateur est $-2$, et $\eu^{\iu j\pi/n} - 1 =
\eu^{\iu j\pi/2n}\cdot2\iu\sin\frac{j\pi}{2n}$, donc
l'expression entière est $\frac{-2}{2\iu\sin(j\pi/2n)} =
\frac{\iu}{\sin(j\pi/2n)}$~: purement imaginaire, $S_j = 0$
encore. D'où la règle à poids égaux $\frac\pi
n\sum_kf(t_k)$ intègre $T_0$ ($\sum w_i = \pi = \int w$) et
tue $T_1, \dots, T_{n-1}$ exactement comme $\int T_jw = 0$~:
elle est exacte jusqu'au \omterm{def:b3:galois:extension}{degré} $n - 1$. Les poids exacts
jusqu'au \omterm{def:b3:galois:extension}{degré} $n-1$ en des nœuds donnés sont uniques (\omterm{def:b3:topology:basis}{base}
de Lagrange)~: les poids de Gauss sont tous $\frac\pi n$.
Pour $n = 3$~: nœuds $\pm\frac{\sqrt3}2, 0$ et
\[
\int_{-1}^1\frac{f(t)}{\sqrt{1 - t^2}}\,\dd t \approx
\frac\pi3\Bigl[f\Bigl(\tfrac{\sqrt3}2\Bigr) + f(0) +
f\Bigl(-\tfrac{\sqrt3}2\Bigr)\Bigr],
\]
exacte jusqu'au \omterm{def:b3:galois:extension}{degré} $5$.

\textbf{18.} Pour un monique $P$ de \omterm{def:b3:galois:extension}{degré} $n$~: $P = p_n + r$
avec $r \in \R_{n-1}[t]$, et $p_n \perp \R_{n-1}[t]$
(question 1), donc $\norm P^2 = \norm{p_n}^2 + \norm r^2 \geq
\norm{p_n}^2$, avec égalité ssi $r = 0$~: $p_n$ est le
résidu de projection orthogonale de $t^n$ sur
$\R_{n-1}[t]^\perp$, i.e.\ le polynôme monique le plus
proche du sous-espace qu'il doit éviter. Le $2^{1-n}T_n$ de
Chebyshev répond à la même question pour la norme uniforme~:
moindre déviation de zéro, une fois en $L^2(w)$, une fois en
$L^\infty$.

\textbf{19.} Écrire $K_n(x, y) =
\sum_{k=0}^n\frac{p_k(x)p_k(y)}{h_k}$. \omterm{def:b3:topology:basis}{Base} $n = 0$~: $(x -
y)\frac1{h_0} = \frac{p_1(x)\cdot1 - 1\cdot p_1(y)}{h_0}$
puisque $p_1 = t - a_0$. Étape~: en supposant l'identité pour
$n - 1$,
\[
(x - y)\,K_n(x,y) = \frac{p_n(x)p_{n-1}(y) -
p_{n-1}(x)p_n(y)}{h_{n-1}} +
\frac{(x - y)\,p_n(x)p_n(y)}{h_n} ;
\]
substituer $x\,p_n(x) = p_{n+1}(x) + a_np_n(x) +
b_np_{n-1}(x)$ et $y\,p_n(y) = p_{n+1}(y) + a_np_n(y) +
b_np_{n-1}(y)$ dans le second terme~: les contributions $a_n$
s'annulent, et les contributions $b_n = \frac{h_n}{h_{n-1}}$
annulent le terme d'induction~; ce qui survit est
$\frac{p_{n+1}(x)p_n(y) - p_n(x)p_{n+1}(y)}{h_n}$. La forme
confluente suit en faisant $y \to x$ (les deux côtés sont
polynomiaux en $y$).

\textbf{20.} La forme confluente donne $p_{n+1}'p_n -
p_n'p_{n+1} = h_n\sum_{k\leq n}\frac{p_k^2}{h_k} \geq
\frac{h_n}{h_0} > 0$ partout. En une racine $x_0$ de
$p_{n+1}$~: $p_{n+1}'(x_0)\,p_n(x_0) > 0$, donc $p_n(x_0)
\neq 0$ (pas de racines communes). Entre des racines
consécutives $x_0 < x_1$ de $p_{n+1}$ (toutes \omterm{def:b3:groups:simple}{simples}, Partie
I), $p_{n+1}'$ a des signes opposés, donc aussi $p_n$~: une
racine de $p_n$ gît dans chacun des $n$ interstices --- et
cela épuise ses $n$ racines~: entrelacement.

\textbf{21.} En développant $D_n(t) = \det(tI_n - J_n)$ le
long de la dernière ligne~: $D_n = (t - a_{n-1})D_{n-1} -
b_{n-1}D_{n-2}$, avec $D_0 = 1$, $D_1 = t - a_0$~: la
récurrence et les germes du monique $p_n$, donc $D_n = p_n$.
Racines de $p_n$ = valeurs propres de la $J_n$ symétrique~:
réelles, et \omterm{def:b3:groups:simple}{simples} par la question 19 --- la quadrature de
Gauss est la théorie spectrale d'une matrice tridiagonale
déguisée, l'ombre de dimension finie du~\cref{ch:b3:spectral}.

\textbf{22.} Dictionnaire~:
\begin{center}
\begin{tabular}{l|c|c|c}
 & Legendre & Hermite & Chebyshev\\
\hline
intervalle & $\intcc{-1}1$ & $\R$ & $\intcc{-1}1$\\
poids & $1$ & $\eu^{-t^2}$ & $(1-t^2)^{-1/2}$\\
formule & Rodrigues & $(-1)^n\eu^{t^2}
\frac{\dd^n}{\dd t^n}\eu^{-t^2}$ & $\cos(n\arccos t)$\\
norme$^2$ & $\frac2{2n+1}$ & $2^nn!\sqrt\pi$ & $\pi,
\frac\pi2$\\
habitat & quadrature & calcul gaussien & minimax\\
\end{tabular}
\end{center}
(chacun avec sa récurrence à trois termes~: forme générale
pour Legendre, $H_{n+1} = 2tH_n - 2nH_{n-1}$, $T_{n+1} =
2tT_n - T_{n-1}$). La théorie générale a fourni ce qu'aucune
famille isolée ne montre~: réalité et entrelacement des
racines, positivité des poids de quadrature, la \omterm{def:b3:groups:simple}{simple}
existence de la récurrence et de Christoffel--Darboux ---
conséquences de l'orthogonalité seule, uniformes en le
poids.

\textbf{23.} Existence~: l'application linéaire $\R_{2n-1}[t]
\to \R^{2n}$, $P \mapsto (P(t_1), P'(t_1), \dots, P(t_n),
P'(t_n))$, est injective (un $P$ dans le noyau a $n$ racines
doubles et \omterm{def:b3:galois:extension}{degré} $\leq 2n - 1$, donc $P = 0$) entre espaces
de dimension égale $2n$~: bijective. Erreur ponctuelle~: fixer
$t$ non nœud et choisir $K$ pour que $g(s) = f(s) - Hf(s) -
K\,p_n(s)^2$ s'annule en $s = t$. Alors $g$ s'annule aux $n +
1$ points distincts $t, t_1, \dots, t_n$, et $g'$ s'annule
aussi en chaque $t_i$ (à la fois $f - Hf$ et $p_n^2$ y ont
des zéros doubles). Rolle donne $n$ zéros de $g'$ strictement
entre les zéros consécutifs de $g$ --- distincts des nœuds ---
donc $g'$ a $2n$ zéros distincts~; en appliquant Rolle $2n -
1$ fois de plus on produit $\xi_t$ avec $g^{(2n)}(\xi_t) =
0$. Comme $\deg Hf \leq 2n - 1$ et $p_n^2$ est monique de
\omterm{def:b3:galois:extension}{degré} $2n$, $g^{(2n)} = f^{(2n)} - K\,(2n)!$, d'où $K =
f^{(2n)}(\xi_t)/(2n)!$ --- et l'identité est triviale aux
nœuds. Intégration~: $Q_n(f) = Q_n(Hf)$ ($Hf$ coïncide avec
$f$ aux nœuds) et $Q_n(Hf) = \int Hf\,w$ par exactitude
jusqu'au \omterm{def:b3:galois:extension}{degré} $2n - 1$ (question 9), donc l'erreur de
quadrature est $\int(f - Hf)\,w$. Avec $m, M$ les extrema de
$f^{(2n)}$ sur $I$, l'identité ponctuelle serre
\[
\frac{m\,h_n}{(2n)!} \;\leq\; \int_I(f - Hf)\,w
\;\leq\; \frac{M\,h_n}{(2n)!} ,
\]
et le théorème des valeurs intermédiaires appliqué à la
\omterm{def:b3:topology:continuity}{continue} $f^{(2n)}$ livre $\xi$. (Pour Legendre avec $n =
2$~: $h_2 = \int_{-1}^1(t^2 - \frac13)^2\dd t =
\frac8{45}$, donc l'erreur est $f^{(4)}(\xi)/135$.)

\textbf{24.} Le noyau reproduit $\R_{n-1}[t]$~: en développant
$q = \sum_k\frac{\langle p_k, q\rangle}{h_k}p_k$ on obtient
\[
\int_I K_n(t_i, t)\,q(t)\,w(t)\dd t = q(t_i)
\]
pour tout $q$ de \omterm{def:b3:galois:extension}{degré} $\leq n - 1$.
Prendre $q = \ell_i$~: le membre de gauche égale
$\ell_i(t_i) = 1$. Mais $t \mapsto K_n(t_i, t)\,\ell_i(t)$
est un polynôme de \omterm{def:b3:galois:extension}{degré}
$\leq (n - 1) + (n - 1) = 2n - 2$, sur lequel $Q_n$ est exacte
(question 9), et il s'annule en tout nœud $t_j \neq t_i$
(facteur $\ell_i$), donc
\[
1 = \int_I K_n(t_i, t)\,\ell_i(t)\,w(t)\dd t
= w_i\,K_n(t_i, t_i)
= w_i\sum_{k=0}^{n-1}\frac{p_k(t_i)^2}{h_k} .
\]
La somme est $> 0$ (son terme $k = 0$ est $1/h_0 > 0$)~: la
formule annoncée, et la positivité encore. Vérification ($n =
2$, Legendre)~: $p_0 = 1$, $h_0 = 2$, $p_1 = t$, $h_1 =
\frac23$~; en $t_i = \pm\frac1{\sqrt3}$,
\[
K_2(t_i, t_i) = \frac12 + \frac{1/3}{2/3} = 1,
\qquad w_i = 1,
\]
comme trouvé à la question 11.

\textbf{25.} En substituant $t = \cos\theta$, l'intégrale est
$\int_0^\pi\cos^6\theta\,\dd\theta =
\pi\,\frac{5\cdot3\cdot1}{6\cdot4\cdot2} = \frac{5\pi}{16}$
(Wallis, \cref{exo:b3:product:8}). La règle
Chebyshev--Gauss $n = 3$ (question 17) a pour nœuds
$\cos\frac\pi6 = \frac{\sqrt3}2$, $\cos\frac\pi2 = 0$,
$\cos\frac{5\pi}6 = -\frac{\sqrt3}2$ et des poids égaux
$\frac\pi3$~:
\[
Q_3(t^6) = \frac\pi3\Bigl(2\cdot\Bigl(\frac{\sqrt3}2
\Bigr)^{6}\Bigr) = \frac\pi3\cdot\frac{54}{64}
= \frac{9\pi}{32},
\qquad
\frac{5\pi}{16} - \frac{9\pi}{32} = \frac\pi{32} .
\]
Prédiction~: le polynôme orthogonal monique de \omterm{def:b3:galois:extension}{degré} $3$ est
$2^{-2}T_3 = t^3 - \frac34t$, avec $h_3 =
\frac1{16}\norm{T_3}_w^2 = \frac1{16}\cdot\frac\pi2 =
\frac\pi{32}$~; et $f = t^6$ a $f^{(6)} = 720 = 6!$
constante, donc la question 23 donne l'erreur
$\frac{6!}{6!}\,h_3 = \frac\pi{32}$ --- sans dépendance en
$\xi$ restante, la formule est forcée d'être exacte, et elle
l'est.
\end{solution}
